\Chapter{107}

\Verse{1}
{
    \zR{话说贾政进内，见了枢密院各位大臣，又见了各位王爷。}
}
{
    \eR{[English source not present in the checked-in Bencraft/Joly files.]}
}
{
}

\Verse{2}
{
    \zR{北静王道：“今日我 们传你来，有遵旨问你的事。”}
}
{
    \eR{[English source not present in the checked-in Bencraft/Joly files.]}
}
{
}

\Verse{3}
{
    \zR{贾政急忙跪下。}
}
{
    \eR{[English source not present in the checked-in Bencraft/Joly files.]}
}
{
}

\Verse{4}
{
    \zR{众大臣便问道：“你哥哥交通外官、 恃强凌弱、纵儿聚赌、强占良民妻女不遂逼死的事，你都知道么？”}
}
{
    \eR{[English source not present in the checked-in Bencraft/Joly files.]}
}
{
}

\Verse{5}
{
    \zR{贾政回道：“犯 官自从主恩钦点学政任满后，查看赈恤，于上年冬底回家，又蒙堂派工程，后又任 江西粮道，题参回都，仍在工部行走，日夜不敢怠惰。}
}
{
    \eR{[English source not present in the checked-in Bencraft/Joly files.]}
}
{
}

\Verse{6}
{
    \zR{一应家务，并未留心伺察， 实在糊涂。}
}
{
    \eR{[English source not present in the checked-in Bencraft/Joly files.]}
}
{
}

\Verse{7}
{
    \zR{不能管教子侄，这就是辜负圣恩。}
}
{
    \eR{[English source not present in the checked-in Bencraft/Joly files.]}
}
{
}

\Verse{8}
{
    \zR{只求主上重重治罪。”}
}
{
    \eR{[English source not present in the checked-in Bencraft/Joly files.]}
}
{
}

\Verse{9}
{
    \zR{北静王据说转 奏。}
}
{
    \eR{[English source not present in the checked-in Bencraft/Joly files.]}
}
{
}

\Verse{10}
{
    \zR{不多时传出旨来，北静王便述道：“主上因御史参奏贾赦交通外官，恃强凌弱 ――据该御史指出平安州互相往来，贾赦包揽词讼――严鞫贾赦，据供平安州原系
        姻亲来往，并未干涉官事，该御史亦不能指实。}
}
{
    \eR{[English source not present in the checked-in Bencraft/Joly files.]}
}
{
}

\Verse{11}
{
    \zR{惟有倚势强索石呆子古扇一款是实 的，然系玩物，究非强索良民之物可比。}
}
{
    \eR{[English source not present in the checked-in Bencraft/Joly files.]}
}
{
}

\Verse{12}
{
    \zR{虽石呆子自尽，亦系疯傻所致，与逼勒致 死者有间。}
}
{
    \eR{[English source not present in the checked-in Bencraft/Joly files.]}
}
{
}

\Verse{13}
{
    \zR{今从宽将贾赦发往台站效力赎罪。}
}
{
    \eR{[English source not present in the checked-in Bencraft/Joly files.]}
}
{
}

\Verse{14}
{
    \zR{所参贾珍强占良民妻女为妾不从逼死 一款，提取都察院原案，看得尤二姐实系张华指腹为婚未娶之妻，因伊贫苦自愿退 婚，尤二姐之母愿结贾珍之弟为妾，并非强占。}
}
{
    \eR{[English source not present in the checked-in Bencraft/Joly files.]}
}
{
}

\Verse{15}
{
    \zR{再尤三姐自刎掩埋、并未报官一款， 查尤三姐原系贾珍妻妹，本意为伊择配，因被逼索定礼，众人扬言秽乱，以致羞忿 自尽，并非贾珍逼勒致死。}
}
{
    \eR{[English source not present in the checked-in Bencraft/Joly files.]}
}
{
}

\Verse{16}
{
    \zR{但身系世袭职员，罔知法纪，私埋人命，本应重治，念 伊究属功臣后裔，不忍加罪，亦从宽革去世职，派往海疆效力赎罪。}
}
{
    \eR{[English source not present in the checked-in Bencraft/Joly files.]}
}
{
}

\Verse{17}
{
    \zR{贾蓉年幼无干， 省释。}
}
{
    \eR{[English source not present in the checked-in Bencraft/Joly files.]}
}
{
}

\Verse{18}
{
    \zR{贾政实系在外任多年，居官尚属勤慎，免治伊治家不正之罪。”}
}
{
    \eR{[English source not present in the checked-in Bencraft/Joly files.]}
}
{
}

\Verse{19}
{
    \zR{贾政听了，感激涕零，叩首不及，又叩求王爷代奏下忱。}
}
{
    \eR{[English source not present in the checked-in Bencraft/Joly files.]}
}
{
}

\Verse{20}
{
    \zR{北静王道：“你该叩 谢天恩，更有何奏？”}
}
{
    \eR{[English source not present in the checked-in Bencraft/Joly files.]}
}
{
}

\Verse{21}
{
    \zR{贾政道：“犯官仰蒙圣恩，不加大罪，又蒙将家产给还，实 在扪心惶愧。}
}
{
    \eR{[English source not present in the checked-in Bencraft/Joly files.]}
}
{
}

\Verse{22}
{
    \zR{愿将祖宗遗受重禄，积馀置产，一并交官。”}
}
{
    \eR{[English source not present in the checked-in Bencraft/Joly files.]}
}
{
}

\Verse{23}
{
    \zR{北静王道：“主上仁慈待 下，明慎用刑，赏罚无差。}
}
{
    \eR{[English source not present in the checked-in Bencraft/Joly files.]}
}
{
}

\Verse{24}
{
    \zR{如今既蒙莫大深恩，给还财产，你又何必多此一奏？”}
}
{
    \eR{[English source not present in the checked-in Bencraft/Joly files.]}
}
{
}

\Verse{25}
{
    \zR{众官也说不必。}
}
{
    \eR{[English source not present in the checked-in Bencraft/Joly files.]}
}
{
}

\Verse{26}
{
    \zR{贾政便谢了恩，叩谢了王爷出来，恐贾母不放心，急忙赶回。}
}
{
    \eR{[English source not present in the checked-in Bencraft/Joly files.]}
}
{
}

\Verse{27}
{
    \zR{上下 男女人等不知传进贾政是何吉凶，都在外头打听，一见贾政回家，都略略的放心， 也不敢问。}
}
{
    \eR{[English source not present in the checked-in Bencraft/Joly files.]}
}
{
}

\Verse{28}
{
    \zR{只见贾政忙忙的走到贾母跟前，将蒙圣恩宽免的事细细告诉了一遍。}
}
{
    \eR{[English source not present in the checked-in Bencraft/Joly files.]}
}
{
}

\Verse{29}
{
    \zR{贾母虽则 放心，只是两个世职革去，贾赦又往台站效力，贾珍又往海疆，不免又悲伤起来。}
}
{
    \eR{[English source not present in the checked-in Bencraft/Joly files.]}
}
{
}

\Verse{30}
{
    \zR{邢夫人尤氏听见这话，更哭起来。}
}
{
    \eR{[English source not present in the checked-in Bencraft/Joly files.]}
}
{
}

\Verse{31}
{
    \zR{贾政便道：“老太太放心。}
}
{
    \eR{[English source not present in the checked-in Bencraft/Joly files.]}
}
{
}

\Verse{32}
{
    \zR{大哥虽则台站效力， 也是为国家办事，不致受苦，只要办得妥当，就可复职。}
}
{
    \eR{[English source not present in the checked-in Bencraft/Joly files.]}
}
{
}

\Verse{33}
{
    \zR{珍儿正是年轻，很该出力。}
}
{
    \eR{[English source not present in the checked-in Bencraft/Joly files.]}
}
{
}

\Verse{34}
{
    \zR{若不是这样，便是祖父的馀德亦不能久享。”}
}
{
    \eR{[English source not present in the checked-in Bencraft/Joly files.]}
}
{
}

\Verse{35}
{
    \zR{说了些宽慰的话。}
}
{
    \eR{[English source not present in the checked-in Bencraft/Joly files.]}
}
{
}

\Verse{36}
{
    \zR{贾母素来本不大喜 欢贾赦，那边东府贾珍究竟隔了一层，只有邢夫人尤氏痛哭不止。}
}
{
    \eR{[English source not present in the checked-in Bencraft/Joly files.]}
}
{
}

\Verse{37}
{
    \zR{邢夫人想着：“家 产一空，丈夫年老远出，膝下虽有琏儿，又是素来顺他二叔的，如今都靠着二叔， 他两口子自然更顺着那边去了。}
}
{
    \eR{[English source not present in the checked-in Bencraft/Joly files.]}
}
{
}

\Verse{38}
{
    \zR{独我一人孤苦伶仃，怎么好？”}
}
{
    \eR{[English source not present in the checked-in Bencraft/Joly files.]}
}
{
}

\Verse{39}
{
    \zR{那尤氏本来独掌宁 府的家计，除了贾珍，也算是惟他为尊，又与贾珍夫妻相和；如今犯事远出，家财 抄尽，依住荣府，虽则老太太疼爱，终是依人门下。}
}
{
    \eR{[English source not present in the checked-in Bencraft/Joly files.]}
}
{
}

\Verse{40}
{
    \zR{又兼带着佩凤偕鸾，那蓉儿夫 妇也还不能兴家立业。}
}
{
    \eR{[English source not present in the checked-in Bencraft/Joly files.]}
}
{
}

\Verse{41}
{
    \zR{又想起：“二妹妹三妹妹都是琏二爷闹的，如今他们倒安然 无事，依旧夫妻完聚，只剩我们几个，怎么度日？”}
}
{
    \eR{[English source not present in the checked-in Bencraft/Joly files.]}
}
{
}

\Verse{42}
{
    \zR{想到这里，痛哭起来。}
}
{
    \eR{[English source not present in the checked-in Bencraft/Joly files.]}
}
{
}

\Verse{43}
{
    \zR{贾母不 忍，便问贾政道：“你大哥和珍儿现已定案，可能回家？}
}
{
    \eR{[English source not present in the checked-in Bencraft/Joly files.]}
}
{
}

\Verse{44}
{
    \zR{蓉儿既没他的事，也该放出 来了。”}
}
{
    \eR{[English source not present in the checked-in Bencraft/Joly files.]}
}
{
}

\Verse{45}
{
    \zR{贾政道：“若在定例呢，大哥是不能回家的。}
}
{
    \eR{[English source not present in the checked-in Bencraft/Joly files.]}
}
{
}

\Verse{46}
{
    \zR{我已托人徇个私情，叫我大哥 同着侄儿回家，好置办行装，衙门内业已应了。}
}
{
    \eR{[English source not present in the checked-in Bencraft/Joly files.]}
}
{
}

\Verse{47}
{
    \zR{想来蓉儿同着他爷爷父亲一起出来。}
}
{
    \eR{[English source not present in the checked-in Bencraft/Joly files.]}
}
{
}

\Verse{48}
{
    \zR{只请老太太放心，儿子办去。”}
}
{
    \eR{[English source not present in the checked-in Bencraft/Joly files.]}
}
{
}

\Verse{49}
{
    \zR{贾母又道：“我这几年老的不成人了，总没有问过家事。}
}
{
    \eR{[English source not present in the checked-in Bencraft/Joly files.]}
}
{
}

\Verse{50}
{
    \zR{如今东府里是抄了去 了，房子入官不用说；你大哥那边，琏儿那里，也都抄了。}
}
{
    \eR{[English source not present in the checked-in Bencraft/Joly files.]}
}
{
}

\Verse{51}
{
    \zR{咱们西府里的银库和东 省地土，你知道还剩了多少？}
}
{
    \eR{[English source not present in the checked-in Bencraft/Joly files.]}
}
{
}

\Verse{52}
{
    \zR{他两个起身，也得给他们几千银子才好。”}
}
{
    \eR{[English source not present in the checked-in Bencraft/Joly files.]}
}
{
}

\Verse{53}
{
    \zR{贾政正是没 法，听见贾母一问，心想着：“若是说明，又恐老太太着急；若不说明，不用说将 来，只现在怎样办法呢？”}
}
{
    \eR{[English source not present in the checked-in Bencraft/Joly files.]}
}
{
}

\Verse{54}
{
    \zR{想毕，便回道：“若老太太不问，儿子也不敢说。}
}
{
    \eR{[English source not present in the checked-in Bencraft/Joly files.]}
}
{
}

\Verse{55}
{
    \zR{如今 老太太既问到这里，现在琏儿也在这里，昨日儿子已查了：旧库的银子早已虚空， 不但用尽，外头还有亏空。}
}
{
    \eR{[English source not present in the checked-in Bencraft/Joly files.]}
}
{
}

\Verse{56}
{
    \zR{现今大哥这件事，若不花银托人，虽说主上宽恩，只怕 他们爷儿两个也不大好，就是这项银子尚无打算。}
}
{
    \eR{[English source not present in the checked-in Bencraft/Joly files.]}
}
{
}

\Verse{57}
{
    \zR{东省的地亩，早已寅年吃了卯年 的租儿了，一时也弄不过来，只好尽所有蒙圣恩没有动的衣服首饰折变了，给大哥 和珍儿作盘费罢了。}
}
{
    \eR{[English source not present in the checked-in Bencraft/Joly files.]}
}
{
}

\Verse{58}
{
    \zR{过日的事只可再打算。”}
}
{
    \eR{[English source not present in the checked-in Bencraft/Joly files.]}
}
{
}

\Verse{59}
{
    \zR{贾母听了，又急的眼泪直淌。}
}
{
    \eR{[English source not present in the checked-in Bencraft/Joly files.]}
}
{
}

\Verse{60}
{
    \zR{说道：“怎 样着？}
}
{
    \eR{[English source not present in the checked-in Bencraft/Joly files.]}
}
{
}

\Verse{61}
{
    \zR{咱们家到了这个田地了么？}
}
{
    \eR{[English source not present in the checked-in Bencraft/Joly files.]}
}
{
}

\Verse{62}
{
    \zR{我虽没有经过，我想起我家向日比这里还强十倍， 也是摆了几年虚架子，没有出这样事，已经塌下来了，不消一二年就完了!据你说 起来，咱们竟一两年就不能支了？”}
}
{
    \eR{[English source not present in the checked-in Bencraft/Joly files.]}
}
{
}

\Verse{63}
{
    \zR{贾政道：“若是这两个世俸不动，外头还有些 挪移。}
}
{
    \eR{[English source not present in the checked-in Bencraft/Joly files.]}
}
{
}

\Verse{64}
{
    \zR{如今无可指称，谁肯接济？”}
}
{
    \eR{[English source not present in the checked-in Bencraft/Joly files.]}
}
{
}

\Verse{65}
{
    \zR{说着，也泪流满面，“想起亲戚来，用过我们 的，如今都穷了；没有用过我们的，又不肯照应。}
}
{
    \eR{[English source not present in the checked-in Bencraft/Joly files.]}
}
{
}

\Verse{66}
{
    \zR{昨日儿子也没有细查，只看了家 下的人丁册子，别说上头的钱一无所出，那底下的人也养不起许多。”}
}
{
    \eR{[English source not present in the checked-in Bencraft/Joly files.]}
}
{
}

\Verse{67}
{
    \zR{贾母正在忧虑，只见贾赦、贾珍、贾蓉一齐进来给贾母请安。}
}
{
    \eR{[English source not present in the checked-in Bencraft/Joly files.]}
}
{
}

\Verse{68}
{
    \zR{贾母看这般光景， 一只手拉着贾赦，一只手拉着贾珍，便大哭起来。}
}
{
    \eR{[English source not present in the checked-in Bencraft/Joly files.]}
}
{
}

\Verse{69}
{
    \zR{他两人脸上羞惭，又见贾母哭泣， 都跪在地下哭着说道：“儿孙们不长进，将祖上功勋丢了，又累老太太伤心，儿孙 们是死无葬身之地的了！”}
}
{
    \eR{[English source not present in the checked-in Bencraft/Joly files.]}
}
{
}

\Verse{70}
{
    \zR{满屋中人看这光景，又一齐大哭起来。}
}
{
    \eR{[English source not present in the checked-in Bencraft/Joly files.]}
}
{
}

\Verse{71}
{
    \zR{贾政只得劝解：“倒 先要打算他两个的使用。}
}
{
    \eR{[English source not present in the checked-in Bencraft/Joly files.]}
}
{
}

\Verse{72}
{
    \zR{大约在家只可住得一两日，迟则人家就不依了。”}
}
{
    \eR{[English source not present in the checked-in Bencraft/Joly files.]}
}
{
}

\Verse{73}
{
    \zR{老太太 含悲忍泪的说道：“你两个且各自同你们媳妇们说说话儿去罢。”}
}
{
    \eR{[English source not present in the checked-in Bencraft/Joly files.]}
}
{
}

\Verse{74}
{
    \zR{又吩咐贾政道：“这 件事是不能久待的。}
}
{
    \eR{[English source not present in the checked-in Bencraft/Joly files.]}
}
{
}

\Verse{75}
{
    \zR{想来外面挪移，恐不中用，那时误了钦限，怎么好？}
}
{
    \eR{[English source not present in the checked-in Bencraft/Joly files.]}
}
{
}

\Verse{76}
{
    \zR{只好我替 你们打算罢了。}
}
{
    \eR{[English source not present in the checked-in Bencraft/Joly files.]}
}
{
}

\Verse{77}
{
    \zR{就是家中如此乱糟糟的，也不是常法儿。”}
}
{
    \eR{[English source not present in the checked-in Bencraft/Joly files.]}
}
{
}

\Verse{78}
{
    \zR{一面说着，便叫鸳鸯吩 咐去了。}
}
{
    \eR{[English source not present in the checked-in Bencraft/Joly files.]}
}
{
}

\Verse{79}
{
    \zR{这里贾赦等出来，又与贾政哭泣了一会，都不免将从前任性、过后恼悔、 如今分离的话说了一会，各自夫妻们那边悲伤去了。}
}
{
    \eR{[English source not present in the checked-in Bencraft/Joly files.]}
}
{
}

\Verse{80}
{
    \zR{贾赦年老，倒还撂的下；独有 贾珍与尤氏怎忍分离？}
}
{
    \eR{[English source not present in the checked-in Bencraft/Joly files.]}
}
{
}

\Verse{81}
{
    \zR{贾琏贾蓉两个也只有拉着父亲啼哭。}
}
{
    \eR{[English source not present in the checked-in Bencraft/Joly files.]}
}
{
}

\Verse{82}
{
    \zR{虽说是比军流减等，究 竟生离死别。}
}
{
    \eR{[English source not present in the checked-in Bencraft/Joly files.]}
}
{
}

\Verse{83}
{
    \zR{这也是事到如此，只得大家硬着心肠过去。}
}
{
    \eR{[English source not present in the checked-in Bencraft/Joly files.]}
}
{
}

\Verse{84}
{
    \zR{却说贾母叫邢王二夫人同着鸳鸯等开箱倒笼，将做媳妇到如今积攒的东西都拿 出来，又叫贾赦、贾政、贾珍等一一的分派。}
}
{
    \eR{[English source not present in the checked-in Bencraft/Joly files.]}
}
{
}

\Verse{85}
{
    \zR{给贾赦三千两，说：“这里现有的银 子你拿二千两去做你的盘费使用，留一千给大太太零用。}
}
{
    \eR{[English source not present in the checked-in Bencraft/Joly files.]}
}
{
}

\Verse{86}
{
    \zR{这三千给珍儿：你只许拿 一千去，留下二千给你媳妇收着。}
}
{
    \eR{[English source not present in the checked-in Bencraft/Joly files.]}
}
{
}

\Verse{87}
{
    \zR{仍旧各自过日子。}
}
{
    \eR{[English source not present in the checked-in Bencraft/Joly files.]}
}
{
}

\Verse{88}
{
    \zR{房子还是一处住，饭食各自吃 罢。}
}
{
    \eR{[English source not present in the checked-in Bencraft/Joly files.]}
}
{
}

\Verse{89}
{
    \zR{四丫头将来的亲事，还是我的事。}
}
{
    \eR{[English source not present in the checked-in Bencraft/Joly files.]}
}
{
}

\Verse{90}
{
    \zR{只可怜凤丫头操了一辈子心，如今弄的精光， 也给他三千两，叫他自己收着，不许叫琏儿用。}
}
{
    \eR{[English source not present in the checked-in Bencraft/Joly files.]}
}
{
}

\Verse{91}
{
    \zR{如今他还病的神昏气短，叫平儿来 拿去。}
}
{
    \eR{[English source not present in the checked-in Bencraft/Joly files.]}
}
{
}

\Verse{92}
{
    \zR{这是你祖父留下的衣裳，还有我少年穿的衣服首饰，如今我也用不着了。}
}
{
    \eR{[English source not present in the checked-in Bencraft/Joly files.]}
}
{
}

\Verse{93}
{
    \zR{男 的呢，叫大老爷、珍儿、琏儿、蓉儿拿去分了。}
}
{
    \eR{[English source not present in the checked-in Bencraft/Joly files.]}
}
{
}

\Verse{94}
{
    \zR{女的呢，叫大太太、珍儿媳妇、凤 丫头拿了分去。}
}
{
    \eR{[English source not present in the checked-in Bencraft/Joly files.]}
}
{
}

\Verse{95}
{
    \zR{这五百两银子交给琏儿，明年将林丫头的棺材送回南去。”}
}
{
    \eR{[English source not present in the checked-in Bencraft/Joly files.]}
}
{
}

\Verse{96}
{
    \zR{分派定 了，又叫贾政道：“你说外头还该着账呢，这是少不得的，你叫拿这金子变卖偿还。}
}
{
    \eR{[English source not present in the checked-in Bencraft/Joly files.]}
}
{
}

\Verse{97}
{
    \zR{这是他们闹掉了我的。}
}
{
    \eR{[English source not present in the checked-in Bencraft/Joly files.]}
}
{
}

\Verse{98}
{
    \zR{你也是我的儿子，我并不偏向。}
}
{
    \eR{[English source not present in the checked-in Bencraft/Joly files.]}
}
{
}

\Verse{99}
{
    \zR{宝玉已经成了家，我下剩的 这些金银东西，大约还值几千银子，这是都给宝玉的了。}
}
{
    \eR{[English source not present in the checked-in Bencraft/Joly files.]}
}
{
}

\Verse{100}
{
    \zR{珠儿媳妇向来孝顺我，兰 儿也好，我也分给他们些。}
}
{
    \eR{[English source not present in the checked-in Bencraft/Joly files.]}
}
{
}

\Verse{101}
{
    \zR{这就是我的事情完了。”}
}
{
    \eR{[English source not present in the checked-in Bencraft/Joly files.]}
}
{
}

\Verse{102}
{
    \zR{贾政等见母亲如此明断分晰， 俱跪下哭着说：“老太太这么大年纪，儿孙们没点孝顺，承受老祖宗这样恩典，叫 儿孙们更无地自容了。”}
}
{
    \eR{[English source not present in the checked-in Bencraft/Joly files.]}
}
{
}

\Verse{103}
{
    \zR{贾母道：“别瞎说了。}
}
{
    \eR{[English source not present in the checked-in Bencraft/Joly files.]}
}
{
}

\Verse{104}
{
    \zR{要不闹出这个乱儿来，我还收着呢。}
}
{
    \eR{[English source not present in the checked-in Bencraft/Joly files.]}
}
{
}

\Verse{105}
{
    \zR{只是现在家人太多，只有二老爷当差，留几个人就够了。}
}
{
    \eR{[English source not present in the checked-in Bencraft/Joly files.]}
}
{
}

\Verse{106}
{
    \zR{你就吩咐管事的，将人叫 齐了，分派妥当。}
}
{
    \eR{[English source not present in the checked-in Bencraft/Joly files.]}
}
{
}

\Verse{107}
{
    \zR{各家有人就罢了。}
}
{
    \eR{[English source not present in the checked-in Bencraft/Joly files.]}
}
{
}

\Verse{108}
{
    \zR{譬如那时都抄了，怎么样呢？}
}
{
    \eR{[English source not present in the checked-in Bencraft/Joly files.]}
}
{
}

\Verse{109}
{
    \zR{我们里头的，也 要叫人分派，该配人的配人，赏去的赏去。}
}
{
    \eR{[English source not present in the checked-in Bencraft/Joly files.]}
}
{
}

\Verse{110}
{
    \zR{如今虽说这房子不入官，你到底把这园 子交了才是呢。}
}
{
    \eR{[English source not present in the checked-in Bencraft/Joly files.]}
}
{
}

\Verse{111}
{
    \zR{那些地亩还交琏儿清理，该卖的卖，留的留，再不可支架子，做空 头。}
}
{
    \eR{[English source not present in the checked-in Bencraft/Joly files.]}
}
{
}

\Verse{112}
{
    \zR{我索性说了罢：江南甄家还有几两银子，二太太那里收着，该叫人就送去罢。}
}
{
    \eR{[English source not present in the checked-in Bencraft/Joly files.]}
}
{
}

\Verse{113}
{
    \zR{倘或再有点事儿出来，可不是他们‘躲过了风暴又遭了雨’了么？”}
}
{
    \eR{[English source not present in the checked-in Bencraft/Joly files.]}
}
{
}

\Verse{114}
{
    \zR{贾政本是不知 当家立计的人，一听贾母的话，一一领命，心想：“老太太实在真真是理家的人。}
}
{
    \eR{[English source not present in the checked-in Bencraft/Joly files.]}
}
{
}

\Verse{115}
{
    \zR{都是我们这些不长进的闹坏了。”}
}
{
    \eR{[English source not present in the checked-in Bencraft/Joly files.]}
}
{
}

\Verse{116}
{
    \zR{贾政见贾母劳乏，求着老太太歇歇养神。}
}
{
    \eR{[English source not present in the checked-in Bencraft/Joly files.]}
}
{
}

\Verse{117}
{
    \zR{贾母又道：“我所剩的东西也有限， 等我死了，做结果我的使用。}
}
{
    \eR{[English source not present in the checked-in Bencraft/Joly files.]}
}
{
}

\Verse{118}
{
    \zR{下剩的都给伏侍我的丫头。”}
}
{
    \eR{[English source not present in the checked-in Bencraft/Joly files.]}
}
{
}

\Verse{119}
{
    \zR{贾政等听到这里，更加 伤感，大家跪下：“请老太太宽怀。}
}
{
    \eR{[English source not present in the checked-in Bencraft/Joly files.]}
}
{
}

\Verse{120}
{
    \zR{只愿儿子们托老太太的福，过了些时，都邀了 恩眷，那时兢兢业业的治起家来，以赎前愆，奉养老太太到一百岁。”}
}
{
    \eR{[English source not present in the checked-in Bencraft/Joly files.]}
}
{
}

\Verse{121}
{
    \zR{贾母道：“但 愿这样才好，我死了也好见祖宗。}
}
{
    \eR{[English source not present in the checked-in Bencraft/Joly files.]}
}
{
}

\Verse{122}
{
    \zR{你们别打量我是享得富贵受不得贫穷的人哪!不 过这几年看着你们轰轰烈烈，我乐得都不管，说说笑笑，养身子罢了。}
}
{
    \eR{[English source not present in the checked-in Bencraft/Joly files.]}
}
{
}

\Verse{123}
{
    \zR{那知道家运 一败，直到这样!若说外头好看，里头空虚，是我早知道的了，只是‘居移气，养 移体’，一时下不了台就是了。}
}
{
    \eR{[English source not present in the checked-in Bencraft/Joly files.]}
}
{
}

\Verse{124}
{
    \zR{如今借此正好收敛，守住这个门头儿，不然，叫人 笑话。}
}
{
    \eR{[English source not present in the checked-in Bencraft/Joly files.]}
}
{
}

\Verse{125}
{
    \zR{你还不知，只打量我知道穷了，就着急的要死。}
}
{
    \eR{[English source not present in the checked-in Bencraft/Joly files.]}
}
{
}

\Verse{126}
{
    \zR{我心里是想着祖宗莫大的功 勋，无一日不指望你们比祖宗还强，能够守住也罢了。}
}
{
    \eR{[English source not present in the checked-in Bencraft/Joly files.]}
}
{
}

\Verse{127}
{
    \zR{谁知他们爷儿两个做些什么 勾当！”}
}
{
    \eR{[English source not present in the checked-in Bencraft/Joly files.]}
}
{
}

\Verse{128}
{
    \zR{贾母正自长篇大论的说，只见丰儿慌慌张张的跑来回王夫人道：“今早我们奶 奶听见外头的事，哭了一场，如今气都接不上了，平儿叫我来回太太。”}
}
{
    \eR{[English source not present in the checked-in Bencraft/Joly files.]}
}
{
}

\Verse{129}
{
    \zR{丰儿没有 说完，贾母听见，便问：“到底怎么样？”}
}
{
    \eR{[English source not present in the checked-in Bencraft/Joly files.]}
}
{
}

\Verse{130}
{
    \zR{王夫人便代回道：“如今说是不大好。”}
}
{
    \eR{[English source not present in the checked-in Bencraft/Joly files.]}
}
{
}

\Verse{131}
{
    \zR{贾母起身道：“嗳!这些冤家，竟要磨死我了。”}
}
{
    \eR{[English source not present in the checked-in Bencraft/Joly files.]}
}
{
}

\Verse{132}
{
    \zR{说着，叫人扶着，要亲自看去。}
}
{
    \eR{[English source not present in the checked-in Bencraft/Joly files.]}
}
{
}

\Verse{133}
{
    \zR{贾 政急忙拦住劝道：“老太太伤了好一会子心，又分派了好些事，这会子该歇歇儿了。}
}
{
    \eR{[English source not present in the checked-in Bencraft/Joly files.]}
}
{
}

\Verse{134}
{
    \zR{就是孙子媳妇有什么事，叫媳妇瞧去就是了，何必老太太亲身过去呢？}
}
{
    \eR{[English source not present in the checked-in Bencraft/Joly files.]}
}
{
}

\Verse{135}
{
    \zR{倘或再伤感 起来，老太太身上要有一点儿不好，叫做儿子的怎么处呢？”}
}
{
    \eR{[English source not present in the checked-in Bencraft/Joly files.]}
}
{
}

\Verse{136}
{
    \zR{贾母道：“你们各自 出去，等一会子再进来，我还有话说。”}
}
{
    \eR{[English source not present in the checked-in Bencraft/Joly files.]}
}
{
}

\Verse{137}
{
    \zR{贾政不敢多言，只得出来料理兄侄起身的 事，又叫贾琏挑人跟去。}
}
{
    \eR{[English source not present in the checked-in Bencraft/Joly files.]}
}
{
}

\Verse{138}
{
    \zR{这里贾母才叫鸳鸯等派人拿了给凤姐的东西，跟着过来。}
}
{
    \eR{[English source not present in the checked-in Bencraft/Joly files.]}
}
{
}

\Verse{139}
{
    \zR{凤姐正在气厥。}
}
{
    \eR{[English source not present in the checked-in Bencraft/Joly files.]}
}
{
}

\Verse{140}
{
    \zR{平儿哭的眼肿腮红，听见贾母带着王夫人等过来，疾忙出来迎 接。}
}
{
    \eR{[English source not present in the checked-in Bencraft/Joly files.]}
}
{
}

\Verse{141}
{
    \zR{贾母便问：“这会子怎么样了？”}
}
{
    \eR{[English source not present in the checked-in Bencraft/Joly files.]}
}
{
}

\Verse{142}
{
    \zR{平儿恐惊了贾母，便说：“这会子好些儿。”}
}
{
    \eR{[English source not present in the checked-in Bencraft/Joly files.]}
}
{
}

\Verse{143}
{
    \zR{说着，跟了贾母等进来，赶忙先走过来，轻轻的揭开帐子。}
}
{
    \eR{[English source not present in the checked-in Bencraft/Joly files.]}
}
{
}

\Verse{144}
{
    \zR{凤姐开眼瞧着，只见贾 母进来，满心惭愧。}
}
{
    \eR{[English source not present in the checked-in Bencraft/Joly files.]}
}
{
}

\Verse{145}
{
    \zR{先前原打量贾母等恼他，不疼他了，是死活由他的，不料贾母 亲自来瞧，心里一宽，觉那拥塞的气略松动些，便要扎挣坐起。}
}
{
    \eR{[English source not present in the checked-in Bencraft/Joly files.]}
}
{
}

\Verse{146}
{
    \zR{贾母叫平儿按着： “不用动。}
}
{
    \eR{[English source not present in the checked-in Bencraft/Joly files.]}
}
{
}

\Verse{147}
{
    \zR{你好些么？”}
}
{
    \eR{[English source not present in the checked-in Bencraft/Joly files.]}
}
{
}

\Verse{148}
{
    \zR{凤姐含泪道：“我好些了。}
}
{
    \eR{[English source not present in the checked-in Bencraft/Joly files.]}
}
{
}

\Verse{149}
{
    \zR{只是从小儿过来，老太太、太 太怎么样疼我!那知我福气薄，叫神鬼支使的失魂落魄，不能够在老太太、太太跟
        前尽点儿孝心，讨个好儿，还这样把我当人，叫我帮着料理家务，被我闹的七颠八 倒，我还有什么脸见老太太、太太呢？}
}
{
    \eR{[English source not present in the checked-in Bencraft/Joly files.]}
}
{
}

\Verse{150}
{
    \zR{今日老太太、太太亲自过来，我更担不起了， 恐怕该活三天的又折了两天去了。”}
}
{
    \eR{[English source not present in the checked-in Bencraft/Joly files.]}
}
{
}

\Verse{151}
{
    \zR{说着悲咽。}
}
{
    \eR{[English source not present in the checked-in Bencraft/Joly files.]}
}
{
}

\Verse{152}
{
    \zR{贾母道：“那些事原是外头闹起来的， 与你什么相干？}
}
{
    \eR{[English source not present in the checked-in Bencraft/Joly files.]}
}
{
}

\Verse{153}
{
    \zR{就是你的东西被人拿去，这也算不了什么呀。}
}
{
    \eR{[English source not present in the checked-in Bencraft/Joly files.]}
}
{
}

\Verse{154}
{
    \zR{我带了好些东西给你， 你瞧瞧。”}
}
{
    \eR{[English source not present in the checked-in Bencraft/Joly files.]}
}
{
}

\Verse{155}
{
    \zR{说着，叫人拿上来给他瞧。}
}
{
    \eR{[English source not present in the checked-in Bencraft/Joly files.]}
}
{
}

\Verse{156}
{
    \zR{凤姐本是贪得无厌的人，如今被抄净尽，自 然愁苦，又恐人埋怨，正是几不欲生的时候。}
}
{
    \eR{[English source not present in the checked-in Bencraft/Joly files.]}
}
{
}

\Verse{157}
{
    \zR{今见贾母仍旧疼他，王夫人也不嗔怪， 过来安慰他，又想贾琏无事，心下安放好些。}
}
{
    \eR{[English source not present in the checked-in Bencraft/Joly files.]}
}
{
}

\Verse{158}
{
    \zR{便在枕上与贾母磕头，说道：“请老 太太放心。}
}
{
    \eR{[English source not present in the checked-in Bencraft/Joly files.]}
}
{
}

\Verse{159}
{
    \zR{若是我的病托着老太太的福好了，我情愿自己当个粗使的丫头，尽心竭 力的伏侍老太太、太太罢！”}
}
{
    \eR{[English source not present in the checked-in Bencraft/Joly files.]}
}
{
}

\Verse{160}
{
    \zR{贾母听他说的伤心，不免掉下泪来。}
}
{
    \eR{[English source not present in the checked-in Bencraft/Joly files.]}
}
{
}

\Verse{161}
{
    \zR{宝玉是从来没有经过这大风浪的，心下只知安乐、不知忧患的人，如今碰来碰 去，都是哭泣的事，所以他竟比傻子尤甚，见人哭他就哭。}
}
{
    \eR{[English source not present in the checked-in Bencraft/Joly files.]}
}
{
}

\Verse{162}
{
    \zR{凤姐看见众人忧闷，反 倒勉强说几句宽慰贾母的话，求着：“请老太太、太太回去，我略好些过来磕头。”}
}
{
    \eR{[English source not present in the checked-in Bencraft/Joly files.]}
}
{
}

\Verse{163}
{
    \zR{说着，将头仰起。}
}
{
    \eR{[English source not present in the checked-in Bencraft/Joly files.]}
}
{
}

\Verse{164}
{
    \zR{贾母叫平儿：“好生服侍。}
}
{
    \eR{[English source not present in the checked-in Bencraft/Joly files.]}
}
{
}

\Verse{165}
{
    \zR{短什么，到我那里要去。”}
}
{
    \eR{[English source not present in the checked-in Bencraft/Joly files.]}
}
{
}

\Verse{166}
{
    \zR{说着，带了 王夫人将要回到自己房中，只听见两三处哭声。}
}
{
    \eR{[English source not present in the checked-in Bencraft/Joly files.]}
}
{
}

\Verse{167}
{
    \zR{贾母听着，实在不忍，便叫王夫人 散去，叫宝玉：“去见你大爷大哥，送一送就回来。”}
}
{
    \eR{[English source not present in the checked-in Bencraft/Joly files.]}
}
{
}

\Verse{168}
{
    \zR{自己躺在榻上下泪。}
}
{
    \eR{[English source not present in the checked-in Bencraft/Joly files.]}
}
{
}

\Verse{169}
{
    \zR{幸喜鸳鸯 等能用百样言语劝解，贾母暂且安歇。}
}
{
    \eR{[English source not present in the checked-in Bencraft/Joly files.]}
}
{
}

\Verse{170}
{
    \zR{不言贾赦等分离悲痛。}
}
{
    \eR{[English source not present in the checked-in Bencraft/Joly files.]}
}
{
}

\Verse{171}
{
    \zR{那些跟去的人，谁是愿意的？}
}
{
    \eR{[English source not present in the checked-in Bencraft/Joly files.]}
}
{
}

\Verse{172}
{
    \zR{不免心中抱怨，叫苦连天。}
}
{
    \eR{[English source not present in the checked-in Bencraft/Joly files.]}
}
{
}

\Verse{173}
{
    \zR{正是生离果胜死别，看者比受者更加伤心。}
}
{
    \eR{[English source not present in the checked-in Bencraft/Joly files.]}
}
{
}

\Verse{174}
{
    \zR{好好的一个荣国府，闹到人嚎鬼哭。}
}
{
    \eR{[English source not present in the checked-in Bencraft/Joly files.]}
}
{
}

\Verse{175}
{
    \zR{贾 政最循规矩，在伦常上也讲究的，执手分别后，自己先骑马赶至城外，举酒送行， 又叮咛了好些“国家轸恤勋臣，力图报称”的话。}
}
{
    \eR{[English source not present in the checked-in Bencraft/Joly files.]}
}
{
}

\Verse{176}
{
    \zR{贾赦等挥泪分头而别。}
}
{
    \eR{[English source not present in the checked-in Bencraft/Joly files.]}
}
{
}

\Verse{177}
{
    \zR{贾政带了宝玉回家，未及进门，只见门上有好些人在那里乱嚷，说：“今日旨 意：将荣国公世职着贾政承袭。”}
}
{
    \eR{[English source not present in the checked-in Bencraft/Joly files.]}
}
{
}

\Verse{178}
{
    \zR{那些人在那里要喜钱，门上人和他们分争，说：“是 本来的世职，我们本家袭了。}
}
{
    \eR{[English source not present in the checked-in Bencraft/Joly files.]}
}
{
}

\Verse{179}
{
    \zR{有什么喜报？”}
}
{
    \eR{[English source not present in the checked-in Bencraft/Joly files.]}
}
{
}

\Verse{180}
{
    \zR{那些人说道：“那世职的荣耀，比任 什么还难得，你们大老爷闹掉了，想要这个，再不能的了。}
}
{
    \eR{[English source not present in the checked-in Bencraft/Joly files.]}
}
{
}

\Verse{181}
{
    \zR{如今圣人的恩典比天还 大，又赏给二老爷了，这是千载难逢的，怎么不给喜钱？”}
}
{
    \eR{[English source not present in the checked-in Bencraft/Joly files.]}
}
{
}

\Verse{182}
{
    \zR{正闹着，贾政回家，门 上回了。}
}
{
    \eR{[English source not present in the checked-in Bencraft/Joly files.]}
}
{
}

\Verse{183}
{
    \zR{虽则喜欢，究竟是哥哥犯事所致，反觉感极涕零，赶着进内告诉贾母。}
}
{
    \eR{[English source not present in the checked-in Bencraft/Joly files.]}
}
{
}

\Verse{184}
{
    \zR{贾 母自然喜欢，拉着说了些勤黾报恩的话。}
}
{
    \eR{[English source not present in the checked-in Bencraft/Joly files.]}
}
{
}

\Verse{185}
{
    \zR{王夫人正恐贾母伤心，过来安慰，听得世 职复还，也是欢喜。}
}
{
    \eR{[English source not present in the checked-in Bencraft/Joly files.]}
}
{
}

\Verse{186}
{
    \zR{独有邢夫人尤氏心下悲苦，只不好露出来。}
}
{
    \eR{[English source not present in the checked-in Bencraft/Joly files.]}
}
{
}

\Verse{187}
{
    \zR{且说外面这些趋炎奉势的亲戚朋友，先前贾宅有事，都远避不来；今儿贾政袭 职，知圣眷尚好，大家都来贺喜。}
}
{
    \eR{[English source not present in the checked-in Bencraft/Joly files.]}
}
{
}

\Verse{188}
{
    \zR{那知贾政纯厚性成，因他袭哥哥的职，心内反生 烦恼，只知感激天恩。}
}
{
    \eR{[English source not present in the checked-in Bencraft/Joly files.]}
}
{
}

\Verse{189}
{
    \zR{于第二日进内谢恩，到底将赏还府第园子备折奏请入官。}
}
{
    \eR{[English source not present in the checked-in Bencraft/Joly files.]}
}
{
}

\Verse{190}
{
    \zR{内 廷降旨不必，贾政才得放心回家，以后循分供职。}
}
{
    \eR{[English source not present in the checked-in Bencraft/Joly files.]}
}
{
}

\Verse{191}
{
    \zR{但是家计萧条，入不敷出。}
}
{
    \eR{[English source not present in the checked-in Bencraft/Joly files.]}
}
{
}

\Verse{192}
{
    \zR{贾政又不能在外应酬。}
}
{
    \eR{[English source not present in the checked-in Bencraft/Joly files.]}
}
{
}

\Verse{193}
{
    \zR{家人们见贾政忠厚，凤姐抱 病不能理家，贾琏的亏空一日重似一日，难免典房卖地。}
}
{
    \eR{[English source not present in the checked-in Bencraft/Joly files.]}
}
{
}

\Verse{194}
{
    \zR{府内家人几个有钱的，怕 贾琏缠扰，都装穷躲事，甚至告假不来，各自另寻门路。}
}
{
    \eR{[English source not present in the checked-in Bencraft/Joly files.]}
}
{
}

\Verse{195}
{
    \zR{独有一个包勇，虽是新投 到此，恰遇荣府坏事，他倒有些真心办事，见那些人欺瞒主子，便时常不忿。}
}
{
    \eR{[English source not present in the checked-in Bencraft/Joly files.]}
}
{
}

\Verse{196}
{
    \zR{奈他 是个新来乍到的人，一句话也插不上，他便生气，每日吃了就睡。}
}
{
    \eR{[English source not present in the checked-in Bencraft/Joly files.]}
}
{
}

\Verse{197}
{
    \zR{众人嫌他不肯随 和，便在贾政前说他终日贪杯生事，并不当差。}
}
{
    \eR{[English source not present in the checked-in Bencraft/Joly files.]}
}
{
}

\Verse{198}
{
    \zR{贾政道：“随他去罢。}
}
{
    \eR{[English source not present in the checked-in Bencraft/Joly files.]}
}
{
}

\Verse{199}
{
    \zR{原是甄府荐 来，不好意思。}
}
{
    \eR{[English source not present in the checked-in Bencraft/Joly files.]}
}
{
}

\Verse{200}
{
    \zR{横竖家内添这一个人吃饭，虽说穷，也不在他一人身上。”}
}
{
    \eR{[English source not present in the checked-in Bencraft/Joly files.]}
}
{
}

\Verse{201}
{
    \zR{并不叫 驱逐。}
}
{
    \eR{[English source not present in the checked-in Bencraft/Joly files.]}
}
{
}

\Verse{202}
{
    \zR{众人又在贾琏跟前说他怎么样不好，贾琏此时也不敢自作威福，只得由他。}
}
{
    \eR{[English source not present in the checked-in Bencraft/Joly files.]}
}
{
}

\Verse{203}
{
    \zR{忽一日，包勇耐不过，吃了几杯酒，在荣府街上闲逛，见有两个人说话。}
}
{
    \eR{[English source not present in the checked-in Bencraft/Joly files.]}
}
{
}

\Verse{204}
{
    \zR{那人 说道：“你瞧，这么个大府，前儿抄了家，不知如今怎么样了？”}
}
{
    \eR{[English source not present in the checked-in Bencraft/Joly files.]}
}
{
}

\Verse{205}
{
    \zR{那人道：“他家怎 么能败？}
}
{
    \eR{[English source not present in the checked-in Bencraft/Joly files.]}
}
{
}

\Verse{206}
{
    \zR{听见说，里头有位娘娘是他家的姑娘，虽是死了，到底有根基的。}
}
{
    \eR{[English source not present in the checked-in Bencraft/Joly files.]}
}
{
}

\Verse{207}
{
    \zR{况且我 常见他们来往的都是王公侯伯，那里没有照应？}
}
{
    \eR{[English source not present in the checked-in Bencraft/Joly files.]}
}
{
}

\Verse{208}
{
    \zR{就是现在的府尹，前任的兵部，是 他们的一家儿。}
}
{
    \eR{[English source not present in the checked-in Bencraft/Joly files.]}
}
{
}

\Verse{209}
{
    \zR{难道有这些人还护庇不来么？”}
}
{
    \eR{[English source not present in the checked-in Bencraft/Joly files.]}
}
{
}

\Verse{210}
{
    \zR{那人道：“你白住在这里!别人犹可， 独是那个贾大人更了不得。}
}
{
    \eR{[English source not present in the checked-in Bencraft/Joly files.]}
}
{
}

\Verse{211}
{
    \zR{我常见他在两府来往，前儿御史虽参了，主子还叫府尹 查明实迹再办。}
}
{
    \eR{[English source not present in the checked-in Bencraft/Joly files.]}
}
{
}

\Verse{212}
{
    \zR{你说他怎么样？}
}
{
    \eR{[English source not present in the checked-in Bencraft/Joly files.]}
}
{
}

\Verse{213}
{
    \zR{他本沾过两府的好处，怕人说他回护一家儿，他倒 狠狠的踢了一脚，所以两府里才到底抄了。}
}
{
    \eR{[English source not present in the checked-in Bencraft/Joly files.]}
}
{
}

\Verse{214}
{
    \zR{你说如今的世情还了得吗！”}
}
{
    \eR{[English source not present in the checked-in Bencraft/Joly files.]}
}
{
}

\Verse{215}
{
    \zR{两人无心 说闲话，岂知旁边有人跟着听的明白。}
}
{
    \eR{[English source not present in the checked-in Bencraft/Joly files.]}
}
{
}

\Verse{216}
{
    \zR{包勇心下暗想：“天下有这样人!但不知是我 们老爷的什么人？}
}
{
    \eR{[English source not present in the checked-in Bencraft/Joly files.]}
}
{
}

\Verse{217}
{
    \zR{我若见了他，便打他一个死，闹出事来，我承当去。”}
}
{
    \eR{[English source not present in the checked-in Bencraft/Joly files.]}
}
{
}

\Verse{218}
{
    \zR{那包勇正在 酒后胡思乱想，忽听那边喝道而来。}
}
{
    \eR{[English source not present in the checked-in Bencraft/Joly files.]}
}
{
}

\Verse{219}
{
    \zR{包勇远远站着，只见那两人轻轻的说道：“这 来的就是那个贾大人了。”}
}
{
    \eR{[English source not present in the checked-in Bencraft/Joly files.]}
}
{
}

\Verse{220}
{
    \zR{包勇听了，心里怀恨，趁着酒兴，便大声说道：“没良心 的男女!怎么忘了我们贾家的恩了？”}
}
{
    \eR{[English source not present in the checked-in Bencraft/Joly files.]}
}
{
}

\Verse{221}
{
    \zR{雨村在轿内听得一个“贾”字，便留神观看， 见是一个醉汉，也不理会，过去了。}
}
{
    \eR{[English source not present in the checked-in Bencraft/Joly files.]}
}
{
}

\Verse{222}
{
    \zR{那包勇醉着，不知好歹，便得意洋洋回到府中，问起同伴，知是方才见的那位 大人是这府里提拔起来的，“他不念旧恩，反来踢弄咱们家里，见了他骂他几句，
        他竟不敢答言。”}
}
{
    \eR{[English source not present in the checked-in Bencraft/Joly files.]}
}
{
}

\Verse{223}
{
    \zR{那荣府的人本嫌包勇，只是主人不计较他，如今他又在外头惹祸， 正好趁着贾政无事，便将包勇喝酒闹事的话回了贾政。}
}
{
    \eR{[English source not present in the checked-in Bencraft/Joly files.]}
}
{
}

\Verse{224}
{
    \zR{贾政此时正怕风波，听见家 人回禀，便一时生气，叫进包勇来数骂了几句，也不好深沉责罚他，便派去看园， 不许他在外行走。}
}
{
    \eR{[English source not present in the checked-in Bencraft/Joly files.]}
}
{
}

\Verse{225}
{
    \zR{那包勇本是个直爽的脾气，投了主子，他便赤心护主，那知贾政 反倒听了别人的话骂他。}
}
{
    \eR{[English source not present in the checked-in Bencraft/Joly files.]}
}
{
}

\Verse{226}
{
    \zR{他也不敢再辩，只得收拾行李往园中看守浇灌去了。}
}
{
    \eR{[English source not present in the checked-in Bencraft/Joly files.]}
}
{
}

\Verse{227}
{
    \zR{未知后事如何，且听下回分解。}
}
{
    \eR{[English source not present in the checked-in Bencraft/Joly files.]}
}
{
}

\Verse{228}
{
    \zR{(本章完)}
}
{
    \eR{[English source not present in the checked-in Bencraft/Joly files.]}
}
{
}
